\subsection*{Висновки до другого розділу}
\addcontentsline{toc}{subsection}{Висновки до другого розділу}

У другому розділі обґрунтовано організаційно-педагогічні засади навчання розробки імерсивних освітніх ресурсів як цілісну, внутрішньо узгоджену педагогічну систему, що інтегрує дидактичні підходи та принципи, структурно-функціональну модель, педагогічні умови й триетапну методику навчання. Наскрізною теоретичною рамкою усього розділу виступають когнітивно-афективна модель імерсивного навчання CAMIL (шість медіаторних факторів: присутність, агентність, втілення, когнітивне навантаження, мотивація, саморегуляція) та модель технолого-педагогічних знань TPACK, що забезпечують конструктивну узгодженість між теоретичними положеннями першого розділу та їх операціоналізацією у методичному забезпеченні підготовки. Зміст розділу розкриває логіку переходу від теоретико-методологічного осмислення проблеми (розділ~1) до практико-орієнтованого розв'язання у межах фахової підготовки майбутніх учителів інформатики.

Обґрунтовано, що навчання розробки ІОР потребує опори на сукупність одинадцяти дидактичних підходів~--- аксіологічного, діяльнісного, інноваційного, інтегративного, когнітивного, компетентнісного, конструктивістського, контекстного, рефлексивного, синергетичного та системного,~--- згрупованих у чотири кластери: ціннісно-світоглядний (аксіологічний, компетентнісний~--- фактори мотивації та саморегуляції CAMIL), процесуально-діяльнісний (діяльнісний, конструктивістський, контекстний, рефлексивний~--- фактори агентності, втілення, присутності), технологічно-інноваційний (інноваційний, когнітивний, синергетичний~--- фактор когнітивного навантаження, TK-компонент TPACK) та системно-структурний (системний, інтегративний~--- TPK-інтеграція TPACK). Побудована матриця відповідності підходів п'яти сутнісним характеристикам ІОР та шести медіаторним факторам CAMIL підтверджує повноту запропонованої сукупності: кожна характеристика ІОР забезпечується щонайменше трьома підходами, кожен фактор CAMIL~--- щонайменше двома. З кластеризованих підходів виведено вісім принципів навчання~--- науковості, інтерактивності, практичної спрямованості, інтеграції, креативності, варіативності, рефлексивності та безперервності,~--- що формують методологічне підґрунтя методики. Практичне втілення запропонованих засад реалізується у навчальному посібнику з WebAR-розробки, методиці інтеграції моделей машинного навчання у WebAR-ресурси та структурі авторського курсу <<Синтетична реальність в освіті>>.

Розроблено структурно-функціональну модель підготовки, що охоплює п'ять взаємозумовлених блоків~--- мотиваційно-цільовий, теоретико-методологічний, змістовий, організаційно-технологічний і рефлексивно-оцінний,~--- кожен із яких співвіднесений із конкретними факторами CAMIL та компонентами TPACK (мотиваційно-цільовий~--- мотивація, саморегуляція, PK/CK; теоретико-методологічний~--- когнітивне навантаження, TK/PK/CK; змістовий~--- втілення, агентність, CK/TK; організаційно-технологічний~--- агентність, присутність, TK/TPK; рефлексивно-оцінний~--- саморегуляція, TPK/TPACK). Послідовність блоків відтворює трифазну педагогічну структуру (теоретична $\to$ практична $\to$ апробаційна), закладену в 6-компонентній моделі проєктування ІОР. Організаційною основою практичної діяльності обрано цикл досвідного навчання Д.~Колба (конкретний досвід $\to$ рефлексивне спостереження $\to$ абстрактна концептуалізація $\to$ активне експериментування), що реалізується у форматі: лабораторна робота $\to$ рефлексія $\to$ теоретичне узагальнення $\to$ проєктне створення. Готовність майбутніх учителів інформатики до розробки ІОР визначено як інтегративне професійно-особистісне утворення, що охоплює шість компонентів~--- мотиваційно-ціннісний, когнітивний, діяльнісно-проєктувальний, технологічний, інтерактивний та рефлексивно-оцінний. Конкретною реалізацією моделі виступає авторський курс <<Синтетична реальність в освіті>> (18~тижневих модулів) на основі навчального посібника з WebAR.

Визначено чотири педагогічні умови результативності навчання~--- психолого-педагогічну (стимулювання професійної мотивації через діяльність зі створення та педагогічного застосування імерсивних ресурсів), методичну (упровадження цілісної поетапної методики формування умінь проєктувати й педагогічно обґрунтовувати ІОР), технологічну (забезпечення опанування повного циклу WebAR-розробки від базових інструментів до ML-інтеграції) та організаційну (проєктування освітнього процесу на засадах поетапності, міждисциплінарної інтеграції та конструктивної узгодженості цілей, змісту, методів і оцінювання). Кожна умова співвіднесена з факторами CAMIL, компонентами TPACK та компонентами готовності, що адресує ключову прогалину, виявлену в міжнародних систематичних оглядах XR-програм підготовки вчителів,~--- відсутність теоретично обґрунтованого дизайну підготовки. Для операціоналізації педагогічних умов розроблено триетапну методику навчання розробки ІОР: теоретико-аналітичний етап (опанування засад ІОР, порівняльний аналіз інструментів WebAR; фази <<аналіз потреб>> та <<формулювання вимог>> 6-компонентної моделі), практико-розробницький етап (прогресивна послідовність WebAR-проєктів~--- відстеження зображень $\to$ облич $\to$ поверхонь $\to$ ML-інтеграція~--- за циклом Колба; фази <<вибір ПЗ>>, <<планування>>, <<розробка>>) та інтегративно-рефлексивний етап (проєктування ІОР-інтегрованих уроків шкільної інформатики, педагогічна аргументація, взаємне рецензування; фази <<впровадження>> та <<оцінювання>>). Методика узгоджена з п'ятьма першими принципами навчання М.~Меріла (проблемно-центрованість, активізація, демонстрація, застосування, інтеграція); конкретною її реалізацією виступає курс <<Синтетична реальність в освіті>> (18~модулів) на базі навчального посібника з WebAR.

Результати другого розділу засвідчують, що дидактичні підходи та принципи, структурно-функціональна модель, педагогічні умови й триетапна методика навчання утворюють цілісну педагогічну систему, в якій усі елементи перебувають у взаємозумовленій єдності, забезпеченій наскрізним узгодженням через фактори CAMIL та компоненти TPACK. Отримані теоретико-методичні результати створюють підґрунтя для експериментальної перевірки ефективності розробленої методики, обґрунтування критеріїв і рівнів сформованості готовності майбутніх учителів інформатики до розробки ІОР, що становить зміст наступного розділу дослідження.

